\documentclass[paper=a4]{scrartcl}
\usepackage[utf8]{inputenc}
\usepackage{graphicx}
\usepackage[absolute,overlay]{textpos}
\usepackage[ngerman]{babel}
\usepackage[normalem]{ulem}
\usepackage[top=2cm,bottom=2cm,left=3cm,right=3cm]{geometry}
\usepackage{array}
\title{Antrag auf Nutzerzugang zum Medusa-Cluster}
\date{}

\begin{document}
\maketitle
\thispagestyle{empty}
\begin{textblock*}{\linewidth}[0,0](10mm,10mm)%
  \includegraphics[width=\linewidth]{unibrief_logo_nat.png}%
\end{textblock*}

\vspace{-2.5cm}
\setlength{\extrarowheight}{1cm}

\begin{tabular}[]{rl}
  Account-Name & \uline{\hspace{10cm}} \\
  Vorname & \uline{\hspace{10cm}} \\
  Nachname & \uline{\hspace{10cm}} \\
  (Dienst-)E-Mail & \uline{\hspace{10cm}} \\
  Arbeitsgruppe/Projekt & \uline{\hspace{10cm}} \\
  Betreuer/Vorgesetzter & \uline{\hspace{10cm}} \\
\end{tabular}

\vspace{1cm}
\noindent Mit meiner Unterschrift bestätige ich, dass ich über folgende
Sachverhalte belehrt wurde:
\begin{itemize}
  \item Der Cluster wird von mehreren Arbeitsgruppen genutzt und die verfügbaren
    Ressourcen werden von allen Nutzer geteilt. Daher müssen nicht
    benötige Ressourcen zügig frei gegeben werden, um Speicherplatz und
    Rechenzeit für andere Nutzer zugänglich zu machen.
  \item Auf dem Cluster sind gegebenenfalls Probanden-Daten gespeichert, die
    unbedingt vertraulich behandelt werden müssen. Nutzer verpflichten sich,
    ausschließlich auf Daten zuzugreifen, die für ihre eigene Arbeit relevant
    und freigegeben sind, sowie keine Daten ohne explizite Erlaubnis auf externe
    Datenträger zu übertragen.
  \item Nutzer erhalten einen Zugang zum Cluster, der von allen an das Internet
    angeschlossenen Computern aus benutzt werden kann. Alle Nutzer sind
    verpflichtet für diesen Zugang ein hinreichend sicheres Passwort zu
    verwenden (z.B. zufällige Kombination aus Buchstaben, Zahlen und
    Sonderzeichen), sowie die Zugangsdaten keinen anderen Personen zugänglich zu
    machen. Verdacht auf Missbrauch durch Verlust oder Diebstahl der
    Zugangsdaten ist unverzüglich zu melden.
  \item Nutzer melden unverzüglich Fehlfunktionen oder Störungen des
    Cluster-Betriebes an den zuständigen Administrator beziehungsweise dem
    jeweiligen Betreuer.


\end{itemize}
\vspace{3cm}

\hspace{2cm}Datum\hfill Unterschrift\hspace{2cm}
\end{document}

